\QHeading{第1题　口算。}
\par\noindent


\begin{oralgrid}
\OralItem{1}{3.2+5.7}{8.9}
\OralItem{2}{6.5+2.4}{8.9}
\OralItem{3}{8.9-3.5}{5.4}
\OralItem{4}{7.8-4.1}{3.7}
\OralItem{5}{1.6+7.2}{8.8}
\OralItem{6}{9.7-5.6}{4.1}
\OralItem{7}{4.8+3.5}{8.3}
\OralItem{8}{5.6+2.7}{8.3}
\OralItem{9}{7.3+4.9}{12.2}
\OralItem{10}{9.2-3.8}{5.4}
\OralItem{11}{6.1-2.5}{3.6}
\OralItem{12}{8.5-4.7}{3.8}\end{oralgrid}
\par

\QHeading{口算思路补充（对应上表题号）}

\Answer{上面列出得数，下面按小题给出不同计算思路。计数单位法适用于小数，整数题可比较凑整、拆分或商不变等方法。}
\MethodTitle{第1小题：$3.2+5.7$}

\MethodBlock{口算思路一：按相同计数单位计算}{先看成若干个十分之一，按整数加减，最后还原成小数。}{\begin{aligned}[t]&\phantom{={}}3.2+5.7\\
&=(32+57)\div 10\\
&=89\div 10\\
&=8.9\end{aligned}}

\MethodBlock{口算思路二：凑整再减去多加的数}{把第二个数凑成整数，再减去多加的部分。}{\begin{aligned}[t]&\phantom{={}}3.2+5.7\\
&=3.2+6-0.3\\
&=9.2-0.3\\
&=8.9\end{aligned}}
\MethodTitle{第2小题：$6.5+2.4$}

\MethodBlock{口算思路一：按相同计数单位计算}{先看成若干个十分之一，按整数加减，最后还原成小数。}{\begin{aligned}[t]&\phantom{={}}6.5+2.4\\
&=(65+24)\div 10\\
&=89\div 10\\
&=8.9\end{aligned}}
\MethodTitle{第3小题：$8.9-3.5$}

\MethodBlock{口算思路一：按相同计数单位计算}{先看成若干个十分之一，按整数加减，最后还原成小数。}{\begin{aligned}[t]&\phantom{={}}8.9-3.5\\
&=(89-35)\div 10\\
&=54\div 10\\
&=5.4\end{aligned}}

\MethodBlock{口算思路二：补差法}{从减数先补到一个整数，再补到被减数，两段差相加。}{\begin{aligned}[t]&\phantom{={}}8.9-3.5\\
&=(4-3.5)+(8.9-4)\\
&=0.5+4.9\\
&=5.4\end{aligned}}
\MethodTitle{第4小题：$7.8-4.1$}

\MethodBlock{口算思路一：按相同计数单位计算}{先看成若干个十分之一，按整数加减，最后还原成小数。}{\begin{aligned}[t]&\phantom{={}}7.8-4.1\\
&=(78-41)\div 10\\
&=37\div 10\\
&=3.7\end{aligned}}

\MethodBlock{口算思路二：补差法}{从减数先补到一个整数，再补到被减数，两段差相加。}{\begin{aligned}[t]&\phantom{={}}7.8-4.1\\
&=(5-4.1)+(7.8-5)\\
&=0.9+2.8\\
&=3.7\end{aligned}}
\MethodTitle{第5小题：$1.6+7.2$}

\MethodBlock{口算思路一：按相同计数单位计算}{先看成若干个十分之一，按整数加减，最后还原成小数。}{\begin{aligned}[t]&\phantom{={}}1.6+7.2\\
&=(16+72)\div 10\\
&=88\div 10\\
&=8.8\end{aligned}}
\MethodTitle{第6小题：$9.7-5.6$}

\MethodBlock{口算思路一：按相同计数单位计算}{先看成若干个十分之一，按整数加减，最后还原成小数。}{\begin{aligned}[t]&\phantom{={}}9.7-5.6\\
&=(97-56)\div 10\\
&=41\div 10\\
&=4.1\end{aligned}}

\MethodBlock{口算思路二：补差法}{从减数先补到一个整数，再补到被减数，两段差相加。}{\begin{aligned}[t]&\phantom{={}}9.7-5.6\\
&=(6-5.6)+(9.7-6)\\
&=0.4+3.7\\
&=4.1\end{aligned}}
\MethodTitle{第7小题：$4.8+3.5$}

\MethodBlock{口算思路一：按相同计数单位计算}{先看成若干个十分之一，按整数加减，最后还原成小数。}{\begin{aligned}[t]&\phantom{={}}4.8+3.5\\
&=(48+35)\div 10\\
&=83\div 10\\
&=8.3\end{aligned}}
\MethodTitle{第8小题：$5.6+2.7$}

\MethodBlock{口算思路一：按相同计数单位计算}{先看成若干个十分之一，按整数加减，最后还原成小数。}{\begin{aligned}[t]&\phantom{={}}5.6+2.7\\
&=(56+27)\div 10\\
&=83\div 10\\
&=8.3\end{aligned}}

\MethodBlock{口算思路二：凑整再减去多加的数}{把第二个数凑成整数，再减去多加的部分。}{\begin{aligned}[t]&\phantom{={}}5.6+2.7\\
&=5.6+3-0.3\\
&=8.6-0.3\\
&=8.3\end{aligned}}
\MethodTitle{第9小题：$7.3+4.9$}

\MethodBlock{口算思路一：按相同计数单位计算}{先看成若干个十分之一，按整数加减，最后还原成小数。}{\begin{aligned}[t]&\phantom{={}}7.3+4.9\\
&=(73+49)\div 10\\
&=122\div 10\\
&=12.2\end{aligned}}

\MethodBlock{口算思路二：凑整再减去多加的数}{把第二个数凑成整数，再减去多加的部分。}{\begin{aligned}[t]&\phantom{={}}7.3+4.9\\
&=7.3+5-0.1\\
&=12.3-0.1\\
&=12.2\end{aligned}}
\MethodTitle{第10小题：$9.2-3.8$}

\MethodBlock{口算思路一：按相同计数单位计算}{先看成若干个十分之一，按整数加减，最后还原成小数。}{\begin{aligned}[t]&\phantom{={}}9.2-3.8\\
&=(92-38)\div 10\\
&=54\div 10\\
&=5.4\end{aligned}}

\MethodBlock{口算思路二：补差法}{从减数先补到一个整数，再补到被减数，两段差相加。}{\begin{aligned}[t]&\phantom{={}}9.2-3.8\\
&=(4-3.8)+(9.2-4)\\
&=0.2+5.2\\
&=5.4\end{aligned}}
\MethodTitle{第11小题：$6.1-2.5$}

\MethodBlock{口算思路一：按相同计数单位计算}{先看成若干个十分之一，按整数加减，最后还原成小数。}{\begin{aligned}[t]&\phantom{={}}6.1-2.5\\
&=(61-25)\div 10\\
&=36\div 10\\
&=3.6\end{aligned}}

\MethodBlock{口算思路二：补差法}{从减数先补到一个整数，再补到被减数，两段差相加。}{\begin{aligned}[t]&\phantom{={}}6.1-2.5\\
&=(3-2.5)+(6.1-3)\\
&=0.5+3.1\\
&=3.6\end{aligned}}
\MethodTitle{第12小题：$8.5-4.7$}

\MethodBlock{口算思路一：按相同计数单位计算}{先看成若干个十分之一，按整数加减，最后还原成小数。}{\begin{aligned}[t]&\phantom{={}}8.5-4.7\\
&=(85-47)\div 10\\
&=38\div 10\\
&=3.8\end{aligned}}

\MethodBlock{口算思路二：补差法}{从减数先补到一个整数，再补到被减数，两段差相加。}{\begin{aligned}[t]&\phantom{={}}8.5-4.7\\
&=(5-4.7)+(8.5-5)\\
&=0.3+3.5\\
&=3.8\end{aligned}}


\par\noindent

\QHeading{第2题　列竖式计算，最后一行要验算。}

\Problem{（1）$5.8+2.46=8.26$}
\par


\MethodTitle{方法一：按题意列竖式}
\Answer{\begin{tikzpicture}[x=1em,y=1em,baseline=(current bounding box.center),draw=answerred,text=answerred]
\node[anchor=base,inner sep=0pt] at (0.000,0) {$5$};
\node[anchor=base,inner sep=0pt] at (0.800,0) {$.$};
\node[anchor=base,inner sep=0pt] at (1.100,0) {$8$};
\node[anchor=base,inner sep=0pt] at (1.900,0) {$0$};
\node[anchor=base,inner sep=0pt] at (0.000,-1.1) {$2$};
\node[anchor=base,inner sep=0pt] at (0.800,-1.1) {$.$};
\node[anchor=base,inner sep=0pt] at (1.100,-1.1) {$4$};
\node[anchor=base,inner sep=0pt] at (1.900,-1.1) {$6$};
\node[anchor=base,inner sep=0pt] at (0.000,-2.3) {$8$};
\node[anchor=base,inner sep=0pt] at (0.800,-2.3) {$.$};
\node[anchor=base,inner sep=0pt] at (1.100,-2.3) {$2$};
\node[anchor=base,inner sep=0pt] at (1.900,-2.3) {$6$};
\node[anchor=base,inner sep=0pt] at (-.7,-1.1) {$+$};
\draw[line width=.5pt] (-1,-1.52)--(2.500,-1.52);
\node[anchor=base,inner sep=0pt,font=\normalsize] at (0.000,.72) {$1$};
\end{tikzpicture}}

\MethodBlock{方法二：拆成整数部分和小数部分}{上方竖式及题目要求的验算仍保留；本法用于核对。}{\begin{aligned}[t]&\phantom{={}}5.8+2.46\\
&=5.8+2+0.46\\
&=7.8+0.46\\
&=8.26\end{aligned}}

\MethodBlock{方法三：按相同计数单位计算}{上方竖式及题目要求的验算仍保留；本法用于核对。先看成若干个百分之一，按整数加减，最后还原成小数。}{\begin{aligned}[t]&\phantom{={}}5.8+2.46\\
&=(580+246)\div 100\\
&=826\div 100\\
&=8.26\end{aligned}}
\Problem{（2）$7.8+1.09=8.89$}

\MethodTitle{方法一：按题意列竖式}
\Answer{\begin{tikzpicture}[x=1em,y=1em,baseline=(current bounding box.center),draw=answerred,text=answerred]
\node[anchor=base,inner sep=0pt] at (0.000,0) {$7$};
\node[anchor=base,inner sep=0pt] at (0.800,0) {$.$};
\node[anchor=base,inner sep=0pt] at (1.100,0) {$8$};
\node[anchor=base,inner sep=0pt] at (1.900,0) {$0$};
\node[anchor=base,inner sep=0pt] at (0.000,-1.1) {$1$};
\node[anchor=base,inner sep=0pt] at (0.800,-1.1) {$.$};
\node[anchor=base,inner sep=0pt] at (1.100,-1.1) {$0$};
\node[anchor=base,inner sep=0pt] at (1.900,-1.1) {$9$};
\node[anchor=base,inner sep=0pt] at (0.000,-2.3) {$8$};
\node[anchor=base,inner sep=0pt] at (0.800,-2.3) {$.$};
\node[anchor=base,inner sep=0pt] at (1.100,-2.3) {$8$};
\node[anchor=base,inner sep=0pt] at (1.900,-2.3) {$9$};
\node[anchor=base,inner sep=0pt] at (-.7,-1.1) {$+$};
\draw[line width=.5pt] (-1,-1.52)--(2.500,-1.52);
\end{tikzpicture}}

\MethodBlock{方法二：拆成整数部分和小数部分}{上方竖式及题目要求的验算仍保留；本法用于核对。}{\begin{aligned}[t]&\phantom{={}}7.8+1.09\\
&=7.8+1+0.09\\
&=8.8+0.09\\
&=8.89\end{aligned}}

\MethodBlock{方法三：按相同计数单位计算}{上方竖式及题目要求的验算仍保留；本法用于核对。先看成若干个百分之一，按整数加减，最后还原成小数。}{\begin{aligned}[t]&\phantom{={}}7.8+1.09\\
&=(780+109)\div 100\\
&=889\div 100\\
&=8.89\end{aligned}}
\Problem{（3）$9.5-3.27=6.23$}

\MethodTitle{方法一：按题意列竖式}
\Answer{\begin{tikzpicture}[x=1em,y=1em,baseline=(current bounding box.center),draw=answerred,text=answerred]
\node[anchor=base,inner sep=0pt] at (0.000,0) {$9$};
\node[anchor=base,inner sep=0pt] at (0.800,0) {$.$};
\node[anchor=base,inner sep=0pt] at (1.100,0) {$5$};
\node[anchor=base,inner sep=0pt] at (1.900,0) {$0$};
\node[anchor=base,inner sep=0pt] at (0.000,-1.1) {$3$};
\node[anchor=base,inner sep=0pt] at (0.800,-1.1) {$.$};
\node[anchor=base,inner sep=0pt] at (1.100,-1.1) {$2$};
\node[anchor=base,inner sep=0pt] at (1.900,-1.1) {$7$};
\node[anchor=base,inner sep=0pt] at (0.000,-2.3) {$6$};
\node[anchor=base,inner sep=0pt] at (0.800,-2.3) {$.$};
\node[anchor=base,inner sep=0pt] at (1.100,-2.3) {$2$};
\node[anchor=base,inner sep=0pt] at (1.900,-2.3) {$3$};
\node[anchor=base,inner sep=0pt] at (-.7,-1.1) {$-$};
\draw[line width=.5pt] (-1,-1.52)--(2.500,-1.52);
\draw[line width=.35pt] (0.940,.03)--(1.260,.43);
\node[anchor=base,inner sep=0pt,font=\normalsize] at (1.100,.72) {$4$};
\draw[line width=.35pt] (1.740,.03)--(2.060,.43);
\node[anchor=base,inner sep=0pt,font=\normalsize] at (1.900,.72) {$10$};
\end{tikzpicture}}

\MethodBlock{方法二：拆成整数部分和小数部分}{上方竖式及题目要求的验算仍保留；本法用于核对。}{\begin{aligned}[t]&\phantom{={}}9.5-3.27\\
&=9.5-3-0.27\\
&=6.5-0.27\\
&=6.23\end{aligned}}

\MethodBlock{方法三：按相同计数单位计算}{上方竖式及题目要求的验算仍保留；本法用于核对。先看成若干个百分之一，按整数加减，最后还原成小数。}{\begin{aligned}[t]&\phantom{={}}9.5-3.27\\
&=(950-327)\div 100\\
&=623\div 100\\
&=6.23\end{aligned}}
\Problem{（4）$6.1-2.58=3.52$}

\MethodTitle{方法一：按题意列竖式}
\Answer{\begin{tikzpicture}[x=1em,y=1em,baseline=(current bounding box.center),draw=answerred,text=answerred]
\node[anchor=base,inner sep=0pt] at (0.000,0) {$6$};
\node[anchor=base,inner sep=0pt] at (0.800,0) {$.$};
\node[anchor=base,inner sep=0pt] at (1.100,0) {$1$};
\node[anchor=base,inner sep=0pt] at (1.900,0) {$0$};
\node[anchor=base,inner sep=0pt] at (0.000,-1.1) {$2$};
\node[anchor=base,inner sep=0pt] at (0.800,-1.1) {$.$};
\node[anchor=base,inner sep=0pt] at (1.100,-1.1) {$5$};
\node[anchor=base,inner sep=0pt] at (1.900,-1.1) {$8$};
\node[anchor=base,inner sep=0pt] at (0.000,-2.3) {$3$};
\node[anchor=base,inner sep=0pt] at (0.800,-2.3) {$.$};
\node[anchor=base,inner sep=0pt] at (1.100,-2.3) {$5$};
\node[anchor=base,inner sep=0pt] at (1.900,-2.3) {$2$};
\node[anchor=base,inner sep=0pt] at (-.7,-1.1) {$-$};
\draw[line width=.5pt] (-1,-1.52)--(2.500,-1.52);
\draw[line width=.35pt] (-0.160,.03)--(0.160,.43);
\node[anchor=base,inner sep=0pt,font=\normalsize] at (0.000,.72) {$5$};
\draw[line width=.35pt] (0.940,.03)--(1.260,.43);
\node[anchor=base,inner sep=0pt,font=\normalsize] at (1.100,.72) {$10$};
\draw[line width=.35pt] (1.740,.03)--(2.060,.43);
\node[anchor=base,inner sep=0pt,font=\normalsize] at (1.900,.72) {$10$};
\end{tikzpicture}}

\MethodBlock{方法二：拆成整数部分和小数部分}{上方竖式及题目要求的验算仍保留；本法用于核对。}{\begin{aligned}[t]&\phantom{={}}6.1-2.58\\
&=6.1-2-0.58\\
&=4.1-0.58\\
&=3.52\end{aligned}}

\MethodBlock{方法三：按相同计数单位计算}{上方竖式及题目要求的验算仍保留；本法用于核对。先看成若干个百分之一，按整数加减，最后还原成小数。}{\begin{aligned}[t]&\phantom{={}}6.1-2.58\\
&=(610-258)\div 100\\
&=352\div 100\\
&=3.52\end{aligned}}
\Problem{（5）$3.97+4.65=8.62$}

\MethodTitle{方法一：按题意列竖式}
\Answer{\begin{tikzpicture}[x=1em,y=1em,baseline=(current bounding box.center),draw=answerred,text=answerred]
\node[anchor=base,inner sep=0pt] at (0.000,0) {$3$};
\node[anchor=base,inner sep=0pt] at (0.800,0) {$.$};
\node[anchor=base,inner sep=0pt] at (1.100,0) {$9$};
\node[anchor=base,inner sep=0pt] at (1.900,0) {$7$};
\node[anchor=base,inner sep=0pt] at (0.000,-1.1) {$4$};
\node[anchor=base,inner sep=0pt] at (0.800,-1.1) {$.$};
\node[anchor=base,inner sep=0pt] at (1.100,-1.1) {$6$};
\node[anchor=base,inner sep=0pt] at (1.900,-1.1) {$5$};
\node[anchor=base,inner sep=0pt] at (0.000,-2.3) {$8$};
\node[anchor=base,inner sep=0pt] at (0.800,-2.3) {$.$};
\node[anchor=base,inner sep=0pt] at (1.100,-2.3) {$6$};
\node[anchor=base,inner sep=0pt] at (1.900,-2.3) {$2$};
\node[anchor=base,inner sep=0pt] at (-.7,-1.1) {$+$};
\draw[line width=.5pt] (-1,-1.52)--(2.500,-1.52);
\node[anchor=base,inner sep=0pt,font=\normalsize] at (1.100,.72) {$1$};
\node[anchor=base,inner sep=0pt,font=\normalsize] at (0.000,.72) {$1$};
\end{tikzpicture}}

\MethodBlock{方法二：拆成整数部分和小数部分}{上方竖式及题目要求的验算仍保留；本法用于核对。}{\begin{aligned}[t]&\phantom{={}}3.97+4.65\\
&=3.97+4+0.65\\
&=7.97+0.65\\
&=8.62\end{aligned}}

\MethodBlock{方法三：按相同计数单位计算}{上方竖式及题目要求的验算仍保留；本法用于核对。先看成若干个百分之一，按整数加减，最后还原成小数。}{\begin{aligned}[t]&\phantom{={}}3.97+4.65\\
&=(397+465)\div 100\\
&=862\div 100\\
&=8.62\end{aligned}}
\Problem{（6）$10-2.34=7.66$}

\MethodTitle{方法一：按题意列竖式}
\Answer{\begin{tikzpicture}[x=1em,y=1em,baseline=(current bounding box.center),draw=answerred,text=answerred]
\node[anchor=base,inner sep=0pt] at (0.000,0) {$1$};
\node[anchor=base,inner sep=0pt] at (0.800,0) {$0$};
\node[anchor=base,inner sep=0pt] at (1.600,0) {$.$};
\node[anchor=base,inner sep=0pt] at (1.900,0) {$0$};
\node[anchor=base,inner sep=0pt] at (2.700,0) {$0$};
\node[anchor=base,inner sep=0pt] at (0.800,-1.1) {$2$};
\node[anchor=base,inner sep=0pt] at (1.600,-1.1) {$.$};
\node[anchor=base,inner sep=0pt] at (1.900,-1.1) {$3$};
\node[anchor=base,inner sep=0pt] at (2.700,-1.1) {$4$};
\node[anchor=base,inner sep=0pt] at (0.800,-2.3) {$7$};
\node[anchor=base,inner sep=0pt] at (1.600,-2.3) {$.$};
\node[anchor=base,inner sep=0pt] at (1.900,-2.3) {$6$};
\node[anchor=base,inner sep=0pt] at (2.700,-2.3) {$6$};
\node[anchor=base,inner sep=0pt] at (-.7,-1.1) {$-$};
\draw[line width=.5pt] (-1,-1.52)--(3.300,-1.52);
\draw[line width=.35pt] (-0.160,.03)--(0.160,.43);
\node[anchor=base,inner sep=0pt,font=\normalsize] at (0.000,.72) {$0$};
\draw[line width=.35pt] (0.640,.03)--(0.960,.43);
\node[anchor=base,inner sep=0pt,font=\normalsize] at (0.800,.72) {$9$};
\draw[line width=.35pt] (1.740,.03)--(2.060,.43);
\node[anchor=base,inner sep=0pt,font=\normalsize] at (1.900,.72) {$9$};
\draw[line width=.35pt] (2.540,.03)--(2.860,.43);
\node[anchor=base,inner sep=0pt,font=\normalsize] at (2.700,.72) {$10$};
\end{tikzpicture}}

\MethodBlock{方法二：拆成整数部分和小数部分}{上方竖式及题目要求的验算仍保留；本法用于核对。}{\begin{aligned}[t]&\phantom{={}}10-2.34\\
&=10-2-0.34\\
&=8-0.34\\
&=7.66\end{aligned}}

\MethodBlock{方法三：按相同计数单位计算}{上方竖式及题目要求的验算仍保留；本法用于核对。先看成若干个百分之一，按整数加减，最后还原成小数。}{\begin{aligned}[t]&\phantom{={}}10-2.34\\
&=(1000-234)\div 100\\
&=766\div 100\\
&=7.66\end{aligned}}
\Problem{（7）$4.28+6.5=10.78$}

\MethodTitle{方法一：按题意列竖式}
\Answer{\begin{tikzpicture}[x=1em,y=1em,baseline=(current bounding box.center),draw=answerred,text=answerred]
\node[anchor=base,inner sep=0pt] at (0.800,0) {$4$};
\node[anchor=base,inner sep=0pt] at (1.600,0) {$.$};
\node[anchor=base,inner sep=0pt] at (1.900,0) {$2$};
\node[anchor=base,inner sep=0pt] at (2.700,0) {$8$};
\node[anchor=base,inner sep=0pt] at (0.800,-1.1) {$6$};
\node[anchor=base,inner sep=0pt] at (1.600,-1.1) {$.$};
\node[anchor=base,inner sep=0pt] at (1.900,-1.1) {$5$};
\node[anchor=base,inner sep=0pt] at (2.700,-1.1) {$0$};
\node[anchor=base,inner sep=0pt] at (0.000,-2.3) {$1$};
\node[anchor=base,inner sep=0pt] at (0.800,-2.3) {$0$};
\node[anchor=base,inner sep=0pt] at (1.600,-2.3) {$.$};
\node[anchor=base,inner sep=0pt] at (1.900,-2.3) {$7$};
\node[anchor=base,inner sep=0pt] at (2.700,-2.3) {$8$};
\node[anchor=base,inner sep=0pt] at (-.7,-1.1) {$+$};
\draw[line width=.5pt] (-1,-1.52)--(3.300,-1.52);
\node[anchor=base,inner sep=0pt,font=\normalsize] at (0.000,.72) {$1$};
\end{tikzpicture}\hspace{3pt}验算：\begin{tikzpicture}[x=1em,y=1em,baseline=(current bounding box.center),draw=answerred,text=answerred]
\node[anchor=base,inner sep=0pt] at (0.000,0) {$1$};
\node[anchor=base,inner sep=0pt] at (0.800,0) {$0$};
\node[anchor=base,inner sep=0pt] at (1.600,0) {$.$};
\node[anchor=base,inner sep=0pt] at (1.900,0) {$7$};
\node[anchor=base,inner sep=0pt] at (2.700,0) {$8$};
\node[anchor=base,inner sep=0pt] at (0.800,-1.1) {$6$};
\node[anchor=base,inner sep=0pt] at (1.600,-1.1) {$.$};
\node[anchor=base,inner sep=0pt] at (1.900,-1.1) {$5$};
\node[anchor=base,inner sep=0pt] at (2.700,-1.1) {$0$};
\node[anchor=base,inner sep=0pt] at (0.800,-2.3) {$4$};
\node[anchor=base,inner sep=0pt] at (1.600,-2.3) {$.$};
\node[anchor=base,inner sep=0pt] at (1.900,-2.3) {$2$};
\node[anchor=base,inner sep=0pt] at (2.700,-2.3) {$8$};
\node[anchor=base,inner sep=0pt] at (-.7,-1.1) {$-$};
\draw[line width=.5pt] (-1,-1.52)--(3.300,-1.52);
\draw[line width=.35pt] (-0.160,.03)--(0.160,.43);
\node[anchor=base,inner sep=0pt,font=\normalsize] at (0.000,.72) {$0$};
\draw[line width=.35pt] (0.640,.03)--(0.960,.43);
\node[anchor=base,inner sep=0pt,font=\normalsize] at (0.800,.72) {$10$};
\end{tikzpicture}}

\MethodBlock{方法二：拆成整数部分和小数部分}{上方竖式及题目要求的验算仍保留；本法用于核对。}{\begin{aligned}[t]&\phantom{={}}4.28+6.5\\
&=4.28+6+0.5\\
&=10.28+0.5\\
&=10.78\end{aligned}}

\MethodBlock{方法三：按相同计数单位计算}{上方竖式及题目要求的验算仍保留；本法用于核对。先看成若干个百分之一，按整数加减，最后还原成小数。}{\begin{aligned}[t]&\phantom{={}}4.28+6.5\\
&=(428+650)\div 100\\
&=1078\div 100\\
&=10.78\end{aligned}}
\Problem{（8）$2.07+9.95=12.02$}

\MethodTitle{方法一：按题意列竖式}
\Answer{\begin{tikzpicture}[x=1em,y=1em,baseline=(current bounding box.center),draw=answerred,text=answerred]
\node[anchor=base,inner sep=0pt] at (0.800,0) {$2$};
\node[anchor=base,inner sep=0pt] at (1.600,0) {$.$};
\node[anchor=base,inner sep=0pt] at (1.900,0) {$0$};
\node[anchor=base,inner sep=0pt] at (2.700,0) {$7$};
\node[anchor=base,inner sep=0pt] at (0.800,-1.1) {$9$};
\node[anchor=base,inner sep=0pt] at (1.600,-1.1) {$.$};
\node[anchor=base,inner sep=0pt] at (1.900,-1.1) {$9$};
\node[anchor=base,inner sep=0pt] at (2.700,-1.1) {$5$};
\node[anchor=base,inner sep=0pt] at (0.000,-2.3) {$1$};
\node[anchor=base,inner sep=0pt] at (0.800,-2.3) {$2$};
\node[anchor=base,inner sep=0pt] at (1.600,-2.3) {$.$};
\node[anchor=base,inner sep=0pt] at (1.900,-2.3) {$0$};
\node[anchor=base,inner sep=0pt] at (2.700,-2.3) {$2$};
\node[anchor=base,inner sep=0pt] at (-.7,-1.1) {$+$};
\draw[line width=.5pt] (-1,-1.52)--(3.300,-1.52);
\node[anchor=base,inner sep=0pt,font=\normalsize] at (1.900,.72) {$1$};
\node[anchor=base,inner sep=0pt,font=\normalsize] at (0.800,.72) {$1$};
\node[anchor=base,inner sep=0pt,font=\normalsize] at (0.000,.72) {$1$};
\end{tikzpicture}\hspace{3pt}验算：\begin{tikzpicture}[x=1em,y=1em,baseline=(current bounding box.center),draw=answerred,text=answerred]
\node[anchor=base,inner sep=0pt] at (0.000,0) {$1$};
\node[anchor=base,inner sep=0pt] at (0.800,0) {$2$};
\node[anchor=base,inner sep=0pt] at (1.600,0) {$.$};
\node[anchor=base,inner sep=0pt] at (1.900,0) {$0$};
\node[anchor=base,inner sep=0pt] at (2.700,0) {$2$};
\node[anchor=base,inner sep=0pt] at (0.800,-1.1) {$9$};
\node[anchor=base,inner sep=0pt] at (1.600,-1.1) {$.$};
\node[anchor=base,inner sep=0pt] at (1.900,-1.1) {$9$};
\node[anchor=base,inner sep=0pt] at (2.700,-1.1) {$5$};
\node[anchor=base,inner sep=0pt] at (0.800,-2.3) {$2$};
\node[anchor=base,inner sep=0pt] at (1.600,-2.3) {$.$};
\node[anchor=base,inner sep=0pt] at (1.900,-2.3) {$0$};
\node[anchor=base,inner sep=0pt] at (2.700,-2.3) {$7$};
\node[anchor=base,inner sep=0pt] at (-.7,-1.1) {$-$};
\draw[line width=.5pt] (-1,-1.52)--(3.300,-1.52);
\draw[line width=.35pt] (-0.160,.03)--(0.160,.43);
\node[anchor=base,inner sep=0pt,font=\normalsize] at (0.000,.72) {$0$};
\draw[line width=.35pt] (0.640,.03)--(0.960,.43);
\node[anchor=base,inner sep=0pt,font=\normalsize] at (0.800,.72) {$11$};
\draw[line width=.35pt] (1.740,.03)--(2.060,.43);
\node[anchor=base,inner sep=0pt,font=\normalsize] at (1.900,.72) {$9$};
\draw[line width=.35pt] (2.540,.03)--(2.860,.43);
\node[anchor=base,inner sep=0pt,font=\normalsize] at (2.700,.72) {$12$};
\end{tikzpicture}}

\MethodBlock{方法二：凑整再减去多加的数}{上方竖式及题目要求的验算仍保留；本法用于核对。把第二个数凑成整数，再减去多加的部分。}{\begin{aligned}[t]&\phantom{={}}2.07+9.95\\
&=2.07+10-0.05\\
&=12.07-0.05\\
&=12.02\end{aligned}}

\MethodBlock{方法三：拆成整数部分和小数部分}{上方竖式及题目要求的验算仍保留；本法用于核对。}{\begin{aligned}[t]&\phantom{={}}2.07+9.95\\
&=2.07+9+0.95\\
&=11.07+0.95\\
&=12.02\end{aligned}}
\Problem{（9）$15-7.82=7.18$}

\MethodTitle{方法一：按题意列竖式}
\Answer{\begin{tikzpicture}[x=1em,y=1em,baseline=(current bounding box.center),draw=answerred,text=answerred]
\node[anchor=base,inner sep=0pt] at (0.000,0) {$1$};
\node[anchor=base,inner sep=0pt] at (0.800,0) {$5$};
\node[anchor=base,inner sep=0pt] at (1.600,0) {$.$};
\node[anchor=base,inner sep=0pt] at (1.900,0) {$0$};
\node[anchor=base,inner sep=0pt] at (2.700,0) {$0$};
\node[anchor=base,inner sep=0pt] at (0.800,-1.1) {$7$};
\node[anchor=base,inner sep=0pt] at (1.600,-1.1) {$.$};
\node[anchor=base,inner sep=0pt] at (1.900,-1.1) {$8$};
\node[anchor=base,inner sep=0pt] at (2.700,-1.1) {$2$};
\node[anchor=base,inner sep=0pt] at (0.800,-2.3) {$7$};
\node[anchor=base,inner sep=0pt] at (1.600,-2.3) {$.$};
\node[anchor=base,inner sep=0pt] at (1.900,-2.3) {$1$};
\node[anchor=base,inner sep=0pt] at (2.700,-2.3) {$8$};
\node[anchor=base,inner sep=0pt] at (-.7,-1.1) {$-$};
\draw[line width=.5pt] (-1,-1.52)--(3.300,-1.52);
\draw[line width=.35pt] (-0.160,.03)--(0.160,.43);
\node[anchor=base,inner sep=0pt,font=\normalsize] at (0.000,.72) {$0$};
\draw[line width=.35pt] (0.640,.03)--(0.960,.43);
\node[anchor=base,inner sep=0pt,font=\normalsize] at (0.800,.72) {$14$};
\draw[line width=.35pt] (1.740,.03)--(2.060,.43);
\node[anchor=base,inner sep=0pt,font=\normalsize] at (1.900,.72) {$9$};
\draw[line width=.35pt] (2.540,.03)--(2.860,.43);
\node[anchor=base,inner sep=0pt,font=\normalsize] at (2.700,.72) {$10$};
\end{tikzpicture}\hspace{3pt}验算：\begin{tikzpicture}[x=1em,y=1em,baseline=(current bounding box.center),draw=answerred,text=answerred]
\node[anchor=base,inner sep=0pt] at (0.800,0) {$7$};
\node[anchor=base,inner sep=0pt] at (1.600,0) {$.$};
\node[anchor=base,inner sep=0pt] at (1.900,0) {$1$};
\node[anchor=base,inner sep=0pt] at (2.700,0) {$8$};
\node[anchor=base,inner sep=0pt] at (0.800,-1.1) {$7$};
\node[anchor=base,inner sep=0pt] at (1.600,-1.1) {$.$};
\node[anchor=base,inner sep=0pt] at (1.900,-1.1) {$8$};
\node[anchor=base,inner sep=0pt] at (2.700,-1.1) {$2$};
\node[anchor=base,inner sep=0pt] at (0.000,-2.3) {$1$};
\node[anchor=base,inner sep=0pt] at (0.800,-2.3) {$5$};
\node[anchor=base,inner sep=0pt] at (1.600,-2.3) {$.$};
\node[anchor=base,inner sep=0pt] at (1.900,-2.3) {$0$};
\node[anchor=base,inner sep=0pt] at (2.700,-2.3) {$0$};
\node[anchor=base,inner sep=0pt] at (-.7,-1.1) {$+$};
\draw[line width=.5pt] (-1,-1.52)--(3.300,-1.52);
\node[anchor=base,inner sep=0pt,font=\normalsize] at (1.900,.72) {$1$};
\node[anchor=base,inner sep=0pt,font=\normalsize] at (0.800,.72) {$1$};
\node[anchor=base,inner sep=0pt,font=\normalsize] at (0.000,.72) {$1$};
\end{tikzpicture}}

\MethodBlock{方法二：多减再加回}{上方竖式及题目要求的验算仍保留；本法用于核对。减去凑成的整数后，把多减的部分加回来。}{\begin{aligned}[t]&\phantom{={}}15-7.82\\
&=15-8+0.18\\
&=7+0.18\\
&=7.18\end{aligned}}

\MethodBlock{方法三：拆成整数部分和小数部分}{上方竖式及题目要求的验算仍保留；本法用于核对。}{\begin{aligned}[t]&\phantom{={}}15-7.82\\
&=15-7-0.82\\
&=8-0.82\\
&=7.18\end{aligned}}
\par\noindent

\Answer{
\QHeading{计算过程：2. 列竖式计算}

小数点对齐，缺位补零；从右往左逐列计算。下面完整列出进位、退位及各数位的计算过程。

\par\medskip\noindent
\textbf{第1小题：\(5.8+2.46=8.26\)}\par\smallskip
\noindent
\begin{tikzpicture}[x=1em,y=1em,baseline=(current bounding box.center),draw=answerred,text=answerred]
\node[anchor=base,inner sep=0pt] at (0.000,0) {$5$};
\node[anchor=base,inner sep=0pt] at (0.800,0) {$.$};
\node[anchor=base,inner sep=0pt] at (1.100,0) {$8$};
\node[anchor=base,inner sep=0pt] at (1.900,0) {$0$};
\node[anchor=base,inner sep=0pt] at (0.000,-1.1) {$2$};
\node[anchor=base,inner sep=0pt] at (0.800,-1.1) {$.$};
\node[anchor=base,inner sep=0pt] at (1.100,-1.1) {$4$};
\node[anchor=base,inner sep=0pt] at (1.900,-1.1) {$6$};
\node[anchor=base,inner sep=0pt] at (0.000,-2.3) {$8$};
\node[anchor=base,inner sep=0pt] at (0.800,-2.3) {$.$};
\node[anchor=base,inner sep=0pt] at (1.100,-2.3) {$2$};
\node[anchor=base,inner sep=0pt] at (1.900,-2.3) {$6$};
\node[anchor=base,inner sep=0pt] at (-.7,-1.1) {$+$};
\draw[line width=.5pt] (-1,-1.52)--(2.500,-1.52);
\node[anchor=base,inner sep=0pt,font=\fontsize{6.5}{7}\selectfont] at (0.000,.72) {$1$};
\end{tikzpicture}
\par\smallskip{\fontsize{10}{13}\selectfont 上方小字表示从后一位进来的数。\par}
\vspace{0pt}
\begin{enumerate}[leftmargin=1.7em,itemsep=5pt,topsep=0pt]
\item 百分位：\(\begin{aligned}[t]&\phantom{={}}0+6\\&=6\end{aligned}\)\quad 本位写6。
\item 十分位：\(\begin{aligned}[t]&\phantom{={}}8+4\\&=12\end{aligned}\)\quad 本位写2，向前一位进1。
\item 个位：\(\begin{aligned}[t]&\phantom{={}}5+2+1\\&=7+1\\&=8\end{aligned}\)\quad 本位写8。
\end{enumerate}
\par\smallskip
因此，\(\begin{aligned}[t]&\phantom{={}}5.8+2.46\\&=8.26\end{aligned}\)。
\par\medskip

\par\medskip\noindent
\textbf{第2小题：\(7.8+1.09=8.89\)}\par\smallskip
\noindent
\begin{tikzpicture}[x=1em,y=1em,baseline=(current bounding box.center),draw=answerred,text=answerred]
\node[anchor=base,inner sep=0pt] at (0.000,0) {$7$};
\node[anchor=base,inner sep=0pt] at (0.800,0) {$.$};
\node[anchor=base,inner sep=0pt] at (1.100,0) {$8$};
\node[anchor=base,inner sep=0pt] at (1.900,0) {$0$};
\node[anchor=base,inner sep=0pt] at (0.000,-1.1) {$1$};
\node[anchor=base,inner sep=0pt] at (0.800,-1.1) {$.$};
\node[anchor=base,inner sep=0pt] at (1.100,-1.1) {$0$};
\node[anchor=base,inner sep=0pt] at (1.900,-1.1) {$9$};
\node[anchor=base,inner sep=0pt] at (0.000,-2.3) {$8$};
\node[anchor=base,inner sep=0pt] at (0.800,-2.3) {$.$};
\node[anchor=base,inner sep=0pt] at (1.100,-2.3) {$8$};
\node[anchor=base,inner sep=0pt] at (1.900,-2.3) {$9$};
\node[anchor=base,inner sep=0pt] at (-.7,-1.1) {$+$};
\draw[line width=.5pt] (-1,-1.52)--(2.500,-1.52);
\end{tikzpicture}
\par\smallskip{\fontsize{10}{13}\selectfont 上方小字表示从后一位进来的数。\par}
\vspace{0pt}
\begin{enumerate}[leftmargin=1.7em,itemsep=5pt,topsep=0pt]
\item 百分位：\(\begin{aligned}[t]&\phantom{={}}0+9\\&=9\end{aligned}\)\quad 本位写9。
\item 十分位：\(\begin{aligned}[t]&\phantom{={}}8+0\\&=8\end{aligned}\)\quad 本位写8。
\item 个位：\(\begin{aligned}[t]&\phantom{={}}7+1\\&=8\end{aligned}\)\quad 本位写8。
\end{enumerate}
\par\smallskip
因此，\(\begin{aligned}[t]&\phantom{={}}7.8+1.09\\&=8.89\end{aligned}\)。
\par\medskip

\par\medskip\noindent
\textbf{第3小题：\(9.5-3.27=6.23\)}\par\smallskip
\noindent
\begin{tikzpicture}[x=1em,y=1em,baseline=(current bounding box.center),draw=answerred,text=answerred]
\node[anchor=base,inner sep=0pt] at (0.000,0) {$9$};
\node[anchor=base,inner sep=0pt] at (0.800,0) {$.$};
\node[anchor=base,inner sep=0pt] at (1.100,0) {$5$};
\node[anchor=base,inner sep=0pt] at (1.900,0) {$0$};
\node[anchor=base,inner sep=0pt] at (0.000,-1.1) {$3$};
\node[anchor=base,inner sep=0pt] at (0.800,-1.1) {$.$};
\node[anchor=base,inner sep=0pt] at (1.100,-1.1) {$2$};
\node[anchor=base,inner sep=0pt] at (1.900,-1.1) {$7$};
\node[anchor=base,inner sep=0pt] at (0.000,-2.3) {$6$};
\node[anchor=base,inner sep=0pt] at (0.800,-2.3) {$.$};
\node[anchor=base,inner sep=0pt] at (1.100,-2.3) {$2$};
\node[anchor=base,inner sep=0pt] at (1.900,-2.3) {$3$};
\node[anchor=base,inner sep=0pt] at (-.7,-1.1) {$-$};
\draw[line width=.5pt] (-1,-1.52)--(2.500,-1.52);
\draw[line width=.35pt] (0.940,.03)--(1.260,.43);
\node[anchor=base,inner sep=0pt,font=\fontsize{6.5}{7}\selectfont] at (1.100,.72) {$4$};
\draw[line width=.35pt] (1.740,.03)--(2.060,.43);
\node[anchor=base,inner sep=0pt,font=\fontsize{6.5}{7}\selectfont] at (1.900,.72) {$10$};
\end{tikzpicture}
\par\smallskip{\fontsize{10}{13}\selectfont 划线数字是退位前的数，上方小字是退位后参与计算的数。\par}
\vspace{0pt}
\begin{enumerate}[leftmargin=1.7em,itemsep=5pt,topsep=0pt]
\item 十分位退位：\(\begin{aligned}[t]&\phantom{={}}5-1\\&=4\end{aligned}\)\quad 借出1给后一位。
\item 百分位接收退位：\(\begin{aligned}[t]&\phantom{={}}0+10\\&=10\end{aligned}\)\quad 前一位的1变成本位的10。
\item 百分位：\(\begin{aligned}[t]&\phantom{={}}10-7\\&=3\end{aligned}\)\quad 本位写3。
\item 十分位：\(\begin{aligned}[t]&\phantom{={}}4-2\\&=2\end{aligned}\)\quad 本位写2。
\item 个位：\(\begin{aligned}[t]&\phantom{={}}9-3\\&=6\end{aligned}\)\quad 本位写6。
\end{enumerate}
\par\smallskip
因此，\(\begin{aligned}[t]&\phantom{={}}9.5-3.27\\&=6.23\end{aligned}\)。
\par\medskip

\par\medskip\noindent
\textbf{第4小题：\(6.1-2.58=3.52\)}\par\smallskip
\noindent
\begin{tikzpicture}[x=1em,y=1em,baseline=(current bounding box.center),draw=answerred,text=answerred]
\node[anchor=base,inner sep=0pt] at (0.000,0) {$6$};
\node[anchor=base,inner sep=0pt] at (0.800,0) {$.$};
\node[anchor=base,inner sep=0pt] at (1.100,0) {$1$};
\node[anchor=base,inner sep=0pt] at (1.900,0) {$0$};
\node[anchor=base,inner sep=0pt] at (0.000,-1.1) {$2$};
\node[anchor=base,inner sep=0pt] at (0.800,-1.1) {$.$};
\node[anchor=base,inner sep=0pt] at (1.100,-1.1) {$5$};
\node[anchor=base,inner sep=0pt] at (1.900,-1.1) {$8$};
\node[anchor=base,inner sep=0pt] at (0.000,-2.3) {$3$};
\node[anchor=base,inner sep=0pt] at (0.800,-2.3) {$.$};
\node[anchor=base,inner sep=0pt] at (1.100,-2.3) {$5$};
\node[anchor=base,inner sep=0pt] at (1.900,-2.3) {$2$};
\node[anchor=base,inner sep=0pt] at (-.7,-1.1) {$-$};
\draw[line width=.5pt] (-1,-1.52)--(2.500,-1.52);
\draw[line width=.35pt] (-0.160,.03)--(0.160,.43);
\node[anchor=base,inner sep=0pt,font=\fontsize{6.5}{7}\selectfont] at (0.000,.72) {$5$};
\draw[line width=.35pt] (0.940,.03)--(1.260,.43);
\node[anchor=base,inner sep=0pt,font=\fontsize{6.5}{7}\selectfont] at (1.100,.72) {$10$};
\draw[line width=.35pt] (1.740,.03)--(2.060,.43);
\node[anchor=base,inner sep=0pt,font=\fontsize{6.5}{7}\selectfont] at (1.900,.72) {$10$};
\end{tikzpicture}
\par\smallskip{\fontsize{10}{13}\selectfont 划线数字是退位前的数，上方小字是退位后参与计算的数。\par}
\vspace{0pt}
\begin{enumerate}[leftmargin=1.7em,itemsep=5pt,topsep=0pt]
\item 十分位退位：\(\begin{aligned}[t]&\phantom{={}}1-1\\&=0\end{aligned}\)\quad 借出1给后一位。
\item 百分位接收退位：\(\begin{aligned}[t]&\phantom{={}}0+10\\&=10\end{aligned}\)\quad 前一位的1变成本位的10。
\item 百分位：\(\begin{aligned}[t]&\phantom{={}}10-8\\&=2\end{aligned}\)\quad 本位写2。
\item 个位退位：\(\begin{aligned}[t]&\phantom{={}}6-1\\&=5\end{aligned}\)\quad 借出1给后一位。
\item 十分位接收退位：\(\begin{aligned}[t]&\phantom{={}}0+10\\&=10\end{aligned}\)\quad 前一位的1变成本位的10。
\item 十分位：\(\begin{aligned}[t]&\phantom{={}}10-5\\&=5\end{aligned}\)\quad 本位写5。
\item 个位：\(\begin{aligned}[t]&\phantom{={}}5-2\\&=3\end{aligned}\)\quad 本位写3。
\end{enumerate}
\par\smallskip
因此，\(\begin{aligned}[t]&\phantom{={}}6.1-2.58\\&=3.52\end{aligned}\)。
\par\medskip

\par\medskip\noindent
\textbf{第5小题：\(3.97+4.65=8.62\)}\par\smallskip
\noindent
\begin{tikzpicture}[x=1em,y=1em,baseline=(current bounding box.center),draw=answerred,text=answerred]
\node[anchor=base,inner sep=0pt] at (0.000,0) {$3$};
\node[anchor=base,inner sep=0pt] at (0.800,0) {$.$};
\node[anchor=base,inner sep=0pt] at (1.100,0) {$9$};
\node[anchor=base,inner sep=0pt] at (1.900,0) {$7$};
\node[anchor=base,inner sep=0pt] at (0.000,-1.1) {$4$};
\node[anchor=base,inner sep=0pt] at (0.800,-1.1) {$.$};
\node[anchor=base,inner sep=0pt] at (1.100,-1.1) {$6$};
\node[anchor=base,inner sep=0pt] at (1.900,-1.1) {$5$};
\node[anchor=base,inner sep=0pt] at (0.000,-2.3) {$8$};
\node[anchor=base,inner sep=0pt] at (0.800,-2.3) {$.$};
\node[anchor=base,inner sep=0pt] at (1.100,-2.3) {$6$};
\node[anchor=base,inner sep=0pt] at (1.900,-2.3) {$2$};
\node[anchor=base,inner sep=0pt] at (-.7,-1.1) {$+$};
\draw[line width=.5pt] (-1,-1.52)--(2.500,-1.52);
\node[anchor=base,inner sep=0pt,font=\fontsize{6.5}{7}\selectfont] at (1.100,.72) {$1$};
\node[anchor=base,inner sep=0pt,font=\fontsize{6.5}{7}\selectfont] at (0.000,.72) {$1$};
\end{tikzpicture}
\par\smallskip{\fontsize{10}{13}\selectfont 上方小字表示从后一位进来的数。\par}
\vspace{0pt}
\begin{enumerate}[leftmargin=1.7em,itemsep=5pt,topsep=0pt]
\item 百分位：\(\begin{aligned}[t]&\phantom{={}}7+5\\&=12\end{aligned}\)\quad 本位写2，向前一位进1。
\item 十分位：\(\begin{aligned}[t]&\phantom{={}}9+6+1\\&=15+1\\&=16\end{aligned}\)\quad 本位写6，向前一位进1。
\item 个位：\(\begin{aligned}[t]&\phantom{={}}3+4+1\\&=7+1\\&=8\end{aligned}\)\quad 本位写8。
\end{enumerate}
\par\smallskip
因此，\(\begin{aligned}[t]&\phantom{={}}3.97+4.65\\&=8.62\end{aligned}\)。
\par\medskip

\par\medskip\noindent
\textbf{第6小题：\(10-2.34=7.66\)}\par\smallskip
\noindent
\begin{tikzpicture}[x=1em,y=1em,baseline=(current bounding box.center),draw=answerred,text=answerred]
\node[anchor=base,inner sep=0pt] at (0.000,0) {$1$};
\node[anchor=base,inner sep=0pt] at (0.800,0) {$0$};
\node[anchor=base,inner sep=0pt] at (1.600,0) {$.$};
\node[anchor=base,inner sep=0pt] at (1.900,0) {$0$};
\node[anchor=base,inner sep=0pt] at (2.700,0) {$0$};
\node[anchor=base,inner sep=0pt] at (0.800,-1.1) {$2$};
\node[anchor=base,inner sep=0pt] at (1.600,-1.1) {$.$};
\node[anchor=base,inner sep=0pt] at (1.900,-1.1) {$3$};
\node[anchor=base,inner sep=0pt] at (2.700,-1.1) {$4$};
\node[anchor=base,inner sep=0pt] at (0.800,-2.3) {$7$};
\node[anchor=base,inner sep=0pt] at (1.600,-2.3) {$.$};
\node[anchor=base,inner sep=0pt] at (1.900,-2.3) {$6$};
\node[anchor=base,inner sep=0pt] at (2.700,-2.3) {$6$};
\node[anchor=base,inner sep=0pt] at (-.7,-1.1) {$-$};
\draw[line width=.5pt] (-1,-1.52)--(3.300,-1.52);
\draw[line width=.35pt] (-0.160,.03)--(0.160,.43);
\node[anchor=base,inner sep=0pt,font=\fontsize{6.5}{7}\selectfont] at (0.000,.72) {$0$};
\draw[line width=.35pt] (0.640,.03)--(0.960,.43);
\node[anchor=base,inner sep=0pt,font=\fontsize{6.5}{7}\selectfont] at (0.800,.72) {$9$};
\draw[line width=.35pt] (1.740,.03)--(2.060,.43);
\node[anchor=base,inner sep=0pt,font=\fontsize{6.5}{7}\selectfont] at (1.900,.72) {$9$};
\draw[line width=.35pt] (2.540,.03)--(2.860,.43);
\node[anchor=base,inner sep=0pt,font=\fontsize{6.5}{7}\selectfont] at (2.700,.72) {$10$};
\end{tikzpicture}
\par\smallskip{\fontsize{10}{13}\selectfont 划线数字是退位前的数，上方小字是退位后参与计算的数。\par}
\vspace{0pt}
\begin{enumerate}[leftmargin=1.7em,itemsep=5pt,topsep=0pt]
\item 十位退位：\(\begin{aligned}[t]&\phantom{={}}1-1\\&=0\end{aligned}\)\quad 借出1给后一位。
\item 个位传递退位：\(\begin{aligned}[t]&\phantom{={}}0+10-1\\&=10-1\\&=9\end{aligned}\)\quad 收到10，再借出1给后一位。
\item 十分位传递退位：\(\begin{aligned}[t]&\phantom{={}}0+10-1\\&=10-1\\&=9\end{aligned}\)\quad 收到10，再借出1给后一位。
\item 百分位接收退位：\(\begin{aligned}[t]&\phantom{={}}0+10\\&=10\end{aligned}\)\quad 前一位的1变成本位的10。
\item 百分位：\(\begin{aligned}[t]&\phantom{={}}10-4\\&=6\end{aligned}\)\quad 本位写6。
\item 十分位：\(\begin{aligned}[t]&\phantom{={}}9-3\\&=6\end{aligned}\)\quad 本位写6。
\item 个位：\(\begin{aligned}[t]&\phantom{={}}9-2\\&=7\end{aligned}\)\quad 本位写7。
\item 十位：\(\begin{aligned}[t]&\phantom{={}}0-0\\&=0\end{aligned}\)\quad 本位写0。
\end{enumerate}
\par\smallskip
因此，\(\begin{aligned}[t]&\phantom{={}}10-2.34\\&=7.66\end{aligned}\)。
\par\medskip

\par\medskip\noindent
\textbf{第7小题：\(4.28+6.5=10.78\)}\par\smallskip
\noindent
\begin{tikzpicture}[x=1em,y=1em,baseline=(current bounding box.center),draw=answerred,text=answerred]
\node[anchor=base,inner sep=0pt] at (0.800,0) {$4$};
\node[anchor=base,inner sep=0pt] at (1.600,0) {$.$};
\node[anchor=base,inner sep=0pt] at (1.900,0) {$2$};
\node[anchor=base,inner sep=0pt] at (2.700,0) {$8$};
\node[anchor=base,inner sep=0pt] at (0.800,-1.1) {$6$};
\node[anchor=base,inner sep=0pt] at (1.600,-1.1) {$.$};
\node[anchor=base,inner sep=0pt] at (1.900,-1.1) {$5$};
\node[anchor=base,inner sep=0pt] at (2.700,-1.1) {$0$};
\node[anchor=base,inner sep=0pt] at (0.000,-2.3) {$1$};
\node[anchor=base,inner sep=0pt] at (0.800,-2.3) {$0$};
\node[anchor=base,inner sep=0pt] at (1.600,-2.3) {$.$};
\node[anchor=base,inner sep=0pt] at (1.900,-2.3) {$7$};
\node[anchor=base,inner sep=0pt] at (2.700,-2.3) {$8$};
\node[anchor=base,inner sep=0pt] at (-.7,-1.1) {$+$};
\draw[line width=.5pt] (-1,-1.52)--(3.300,-1.52);
\node[anchor=base,inner sep=0pt,font=\fontsize{6.5}{7}\selectfont] at (0.000,.72) {$1$};
\end{tikzpicture}
\par\smallskip{\fontsize{10}{13}\selectfont 上方小字表示从后一位进来的数。\par}
\vspace{0pt}
\begin{enumerate}[leftmargin=1.7em,itemsep=5pt,topsep=0pt]
\item 百分位：\(\begin{aligned}[t]&\phantom{={}}8+0\\&=8\end{aligned}\)\quad 本位写8。
\item 十分位：\(\begin{aligned}[t]&\phantom{={}}2+5\\&=7\end{aligned}\)\quad 本位写7。
\item 个位：\(\begin{aligned}[t]&\phantom{={}}4+6\\&=10\end{aligned}\)\quad 本位写0，向前一位进1。
\end{enumerate}
最高位再写进上来的1。\par
\par\smallskip
因此，\(\begin{aligned}[t]&\phantom{={}}4.28+6.5\\&=10.78\end{aligned}\)。
\par\medskip

\par\medskip\noindent
\textbf{第7小题验算：\(10.78-6.5\)}\par\smallskip
\noindent
\begin{tikzpicture}[x=1em,y=1em,baseline=(current bounding box.center),draw=answerred,text=answerred]
\node[anchor=base,inner sep=0pt] at (0.000,0) {$1$};
\node[anchor=base,inner sep=0pt] at (0.800,0) {$0$};
\node[anchor=base,inner sep=0pt] at (1.600,0) {$.$};
\node[anchor=base,inner sep=0pt] at (1.900,0) {$7$};
\node[anchor=base,inner sep=0pt] at (2.700,0) {$8$};
\node[anchor=base,inner sep=0pt] at (0.800,-1.1) {$6$};
\node[anchor=base,inner sep=0pt] at (1.600,-1.1) {$.$};
\node[anchor=base,inner sep=0pt] at (1.900,-1.1) {$5$};
\node[anchor=base,inner sep=0pt] at (2.700,-1.1) {$0$};
\node[anchor=base,inner sep=0pt] at (0.800,-2.3) {$4$};
\node[anchor=base,inner sep=0pt] at (1.600,-2.3) {$.$};
\node[anchor=base,inner sep=0pt] at (1.900,-2.3) {$2$};
\node[anchor=base,inner sep=0pt] at (2.700,-2.3) {$8$};
\node[anchor=base,inner sep=0pt] at (-.7,-1.1) {$-$};
\draw[line width=.5pt] (-1,-1.52)--(3.300,-1.52);
\draw[line width=.35pt] (-0.160,.03)--(0.160,.43);
\node[anchor=base,inner sep=0pt,font=\fontsize{6.5}{7}\selectfont] at (0.000,.72) {$0$};
\draw[line width=.35pt] (0.640,.03)--(0.960,.43);
\node[anchor=base,inner sep=0pt,font=\fontsize{6.5}{7}\selectfont] at (0.800,.72) {$10$};
\end{tikzpicture}
\par\smallskip{\fontsize{10}{13}\selectfont 划线数字是退位前的数，上方小字是退位后参与计算的数。\par}
\vspace{0pt}
\begin{enumerate}[leftmargin=1.7em,itemsep=5pt,topsep=0pt]
\item 百分位：\(\begin{aligned}[t]&\phantom{={}}8-0\\&=8\end{aligned}\)\quad 本位写8。
\item 十分位：\(\begin{aligned}[t]&\phantom{={}}7-5\\&=2\end{aligned}\)\quad 本位写2。
\item 十位退位：\(\begin{aligned}[t]&\phantom{={}}1-1\\&=0\end{aligned}\)\quad 借出1给后一位。
\item 个位接收退位：\(\begin{aligned}[t]&\phantom{={}}0+10\\&=10\end{aligned}\)\quad 前一位的1变成本位的10。
\item 个位：\(\begin{aligned}[t]&\phantom{={}}10-6\\&=4\end{aligned}\)\quad 本位写4。
\item 十位：\(\begin{aligned}[t]&\phantom{={}}0-0\\&=0\end{aligned}\)\quad 本位写0。
\end{enumerate}
\par\smallskip
因此，\(\begin{aligned}[t]&\phantom{={}}10.78-6.5\\&=4.28\end{aligned}\)。
\par\medskip

\par\medskip\noindent
\textbf{第8小题：\(2.07+9.95=12.02\)}\par\smallskip
\noindent
\begin{tikzpicture}[x=1em,y=1em,baseline=(current bounding box.center),draw=answerred,text=answerred]
\node[anchor=base,inner sep=0pt] at (0.800,0) {$2$};
\node[anchor=base,inner sep=0pt] at (1.600,0) {$.$};
\node[anchor=base,inner sep=0pt] at (1.900,0) {$0$};
\node[anchor=base,inner sep=0pt] at (2.700,0) {$7$};
\node[anchor=base,inner sep=0pt] at (0.800,-1.1) {$9$};
\node[anchor=base,inner sep=0pt] at (1.600,-1.1) {$.$};
\node[anchor=base,inner sep=0pt] at (1.900,-1.1) {$9$};
\node[anchor=base,inner sep=0pt] at (2.700,-1.1) {$5$};
\node[anchor=base,inner sep=0pt] at (0.000,-2.3) {$1$};
\node[anchor=base,inner sep=0pt] at (0.800,-2.3) {$2$};
\node[anchor=base,inner sep=0pt] at (1.600,-2.3) {$.$};
\node[anchor=base,inner sep=0pt] at (1.900,-2.3) {$0$};
\node[anchor=base,inner sep=0pt] at (2.700,-2.3) {$2$};
\node[anchor=base,inner sep=0pt] at (-.7,-1.1) {$+$};
\draw[line width=.5pt] (-1,-1.52)--(3.300,-1.52);
\node[anchor=base,inner sep=0pt,font=\fontsize{6.5}{7}\selectfont] at (1.900,.72) {$1$};
\node[anchor=base,inner sep=0pt,font=\fontsize{6.5}{7}\selectfont] at (0.800,.72) {$1$};
\node[anchor=base,inner sep=0pt,font=\fontsize{6.5}{7}\selectfont] at (0.000,.72) {$1$};
\end{tikzpicture}
\par\smallskip{\fontsize{10}{13}\selectfont 上方小字表示从后一位进来的数。\par}
\vspace{0pt}
\begin{enumerate}[leftmargin=1.7em,itemsep=5pt,topsep=0pt]
\item 百分位：\(\begin{aligned}[t]&\phantom{={}}7+5\\&=12\end{aligned}\)\quad 本位写2，向前一位进1。
\item 十分位：\(\begin{aligned}[t]&\phantom{={}}0+9+1\\&=9+1\\&=10\end{aligned}\)\quad 本位写0，向前一位进1。
\item 个位：\(\begin{aligned}[t]&\phantom{={}}2+9+1\\&=11+1\\&=12\end{aligned}\)\quad 本位写2，向前一位进1。
\end{enumerate}
最高位再写进上来的1。\par
\par\smallskip
因此，\(\begin{aligned}[t]&\phantom{={}}2.07+9.95\\&=12.02\end{aligned}\)。
\par\medskip

\par\medskip\noindent
\textbf{第8小题验算：\(12.02-9.95\)}\par\smallskip
\noindent
\begin{tikzpicture}[x=1em,y=1em,baseline=(current bounding box.center),draw=answerred,text=answerred]
\node[anchor=base,inner sep=0pt] at (0.000,0) {$1$};
\node[anchor=base,inner sep=0pt] at (0.800,0) {$2$};
\node[anchor=base,inner sep=0pt] at (1.600,0) {$.$};
\node[anchor=base,inner sep=0pt] at (1.900,0) {$0$};
\node[anchor=base,inner sep=0pt] at (2.700,0) {$2$};
\node[anchor=base,inner sep=0pt] at (0.800,-1.1) {$9$};
\node[anchor=base,inner sep=0pt] at (1.600,-1.1) {$.$};
\node[anchor=base,inner sep=0pt] at (1.900,-1.1) {$9$};
\node[anchor=base,inner sep=0pt] at (2.700,-1.1) {$5$};
\node[anchor=base,inner sep=0pt] at (0.800,-2.3) {$2$};
\node[anchor=base,inner sep=0pt] at (1.600,-2.3) {$.$};
\node[anchor=base,inner sep=0pt] at (1.900,-2.3) {$0$};
\node[anchor=base,inner sep=0pt] at (2.700,-2.3) {$7$};
\node[anchor=base,inner sep=0pt] at (-.7,-1.1) {$-$};
\draw[line width=.5pt] (-1,-1.52)--(3.300,-1.52);
\draw[line width=.35pt] (-0.160,.03)--(0.160,.43);
\node[anchor=base,inner sep=0pt,font=\fontsize{6.5}{7}\selectfont] at (0.000,.72) {$0$};
\draw[line width=.35pt] (0.640,.03)--(0.960,.43);
\node[anchor=base,inner sep=0pt,font=\fontsize{6.5}{7}\selectfont] at (0.800,.72) {$11$};
\draw[line width=.35pt] (1.740,.03)--(2.060,.43);
\node[anchor=base,inner sep=0pt,font=\fontsize{6.5}{7}\selectfont] at (1.900,.72) {$9$};
\draw[line width=.35pt] (2.540,.03)--(2.860,.43);
\node[anchor=base,inner sep=0pt,font=\fontsize{6.5}{7}\selectfont] at (2.700,.72) {$12$};
\end{tikzpicture}
\par\smallskip{\fontsize{10}{13}\selectfont 划线数字是退位前的数，上方小字是退位后参与计算的数。\par}
\vspace{0pt}
\begin{enumerate}[leftmargin=1.7em,itemsep=5pt,topsep=0pt]
\item 个位退位：\(\begin{aligned}[t]&\phantom{={}}2-1\\&=1\end{aligned}\)\quad 借出1给后一位。
\item 十分位传递退位：\(\begin{aligned}[t]&\phantom{={}}0+10-1\\&=10-1\\&=9\end{aligned}\)\quad 收到10，再借出1给后一位。
\item 百分位接收退位：\(\begin{aligned}[t]&\phantom{={}}2+10\\&=12\end{aligned}\)\quad 前一位的1变成本位的10。
\item 百分位：\(\begin{aligned}[t]&\phantom{={}}12-5\\&=7\end{aligned}\)\quad 本位写7。
\item 十分位：\(\begin{aligned}[t]&\phantom{={}}9-9\\&=0\end{aligned}\)\quad 本位写0。
\item 十位退位：\(\begin{aligned}[t]&\phantom{={}}1-1\\&=0\end{aligned}\)\quad 借出1给后一位。
\item 个位接收退位：\(\begin{aligned}[t]&\phantom{={}}1+10\\&=11\end{aligned}\)\quad 前一位的1变成本位的10。
\item 个位：\(\begin{aligned}[t]&\phantom{={}}11-9\\&=2\end{aligned}\)\quad 本位写2。
\item 十位：\(\begin{aligned}[t]&\phantom{={}}0-0\\&=0\end{aligned}\)\quad 本位写0。
\end{enumerate}
\par\smallskip
因此，\(\begin{aligned}[t]&\phantom{={}}12.02-9.95\\&=2.07\end{aligned}\)。
\par\medskip

\par\medskip\noindent
\textbf{第9小题：\(15-7.82=7.18\)}\par\smallskip
\noindent
\begin{tikzpicture}[x=1em,y=1em,baseline=(current bounding box.center),draw=answerred,text=answerred]
\node[anchor=base,inner sep=0pt] at (0.000,0) {$1$};
\node[anchor=base,inner sep=0pt] at (0.800,0) {$5$};
\node[anchor=base,inner sep=0pt] at (1.600,0) {$.$};
\node[anchor=base,inner sep=0pt] at (1.900,0) {$0$};
\node[anchor=base,inner sep=0pt] at (2.700,0) {$0$};
\node[anchor=base,inner sep=0pt] at (0.800,-1.1) {$7$};
\node[anchor=base,inner sep=0pt] at (1.600,-1.1) {$.$};
\node[anchor=base,inner sep=0pt] at (1.900,-1.1) {$8$};
\node[anchor=base,inner sep=0pt] at (2.700,-1.1) {$2$};
\node[anchor=base,inner sep=0pt] at (0.800,-2.3) {$7$};
\node[anchor=base,inner sep=0pt] at (1.600,-2.3) {$.$};
\node[anchor=base,inner sep=0pt] at (1.900,-2.3) {$1$};
\node[anchor=base,inner sep=0pt] at (2.700,-2.3) {$8$};
\node[anchor=base,inner sep=0pt] at (-.7,-1.1) {$-$};
\draw[line width=.5pt] (-1,-1.52)--(3.300,-1.52);
\draw[line width=.35pt] (-0.160,.03)--(0.160,.43);
\node[anchor=base,inner sep=0pt,font=\fontsize{6.5}{7}\selectfont] at (0.000,.72) {$0$};
\draw[line width=.35pt] (0.640,.03)--(0.960,.43);
\node[anchor=base,inner sep=0pt,font=\fontsize{6.5}{7}\selectfont] at (0.800,.72) {$14$};
\draw[line width=.35pt] (1.740,.03)--(2.060,.43);
\node[anchor=base,inner sep=0pt,font=\fontsize{6.5}{7}\selectfont] at (1.900,.72) {$9$};
\draw[line width=.35pt] (2.540,.03)--(2.860,.43);
\node[anchor=base,inner sep=0pt,font=\fontsize{6.5}{7}\selectfont] at (2.700,.72) {$10$};
\end{tikzpicture}
\par\smallskip{\fontsize{10}{13}\selectfont 划线数字是退位前的数，上方小字是退位后参与计算的数。\par}
\vspace{0pt}
\begin{enumerate}[leftmargin=1.7em,itemsep=5pt,topsep=0pt]
\item 个位退位：\(\begin{aligned}[t]&\phantom{={}}5-1\\&=4\end{aligned}\)\quad 借出1给后一位。
\item 十分位传递退位：\(\begin{aligned}[t]&\phantom{={}}0+10-1\\&=10-1\\&=9\end{aligned}\)\quad 收到10，再借出1给后一位。
\item 百分位接收退位：\(\begin{aligned}[t]&\phantom{={}}0+10\\&=10\end{aligned}\)\quad 前一位的1变成本位的10。
\item 百分位：\(\begin{aligned}[t]&\phantom{={}}10-2\\&=8\end{aligned}\)\quad 本位写8。
\item 十分位：\(\begin{aligned}[t]&\phantom{={}}9-8\\&=1\end{aligned}\)\quad 本位写1。
\item 十位退位：\(\begin{aligned}[t]&\phantom{={}}1-1\\&=0\end{aligned}\)\quad 借出1给后一位。
\item 个位接收退位：\(\begin{aligned}[t]&\phantom{={}}4+10\\&=14\end{aligned}\)\quad 前一位的1变成本位的10。
\item 个位：\(\begin{aligned}[t]&\phantom{={}}14-7\\&=7\end{aligned}\)\quad 本位写7。
\item 十位：\(\begin{aligned}[t]&\phantom{={}}0-0\\&=0\end{aligned}\)\quad 本位写0。
\end{enumerate}
\par\smallskip
因此，\(\begin{aligned}[t]&\phantom{={}}15-7.82\\&=7.18\end{aligned}\)。
\par\medskip

\par\medskip\noindent
\textbf{第9小题验算：\(7.18+7.82\)}\par\smallskip
\noindent
\begin{tikzpicture}[x=1em,y=1em,baseline=(current bounding box.center),draw=answerred,text=answerred]
\node[anchor=base,inner sep=0pt] at (0.800,0) {$7$};
\node[anchor=base,inner sep=0pt] at (1.600,0) {$.$};
\node[anchor=base,inner sep=0pt] at (1.900,0) {$1$};
\node[anchor=base,inner sep=0pt] at (2.700,0) {$8$};
\node[anchor=base,inner sep=0pt] at (0.800,-1.1) {$7$};
\node[anchor=base,inner sep=0pt] at (1.600,-1.1) {$.$};
\node[anchor=base,inner sep=0pt] at (1.900,-1.1) {$8$};
\node[anchor=base,inner sep=0pt] at (2.700,-1.1) {$2$};
\node[anchor=base,inner sep=0pt] at (0.000,-2.3) {$1$};
\node[anchor=base,inner sep=0pt] at (0.800,-2.3) {$5$};
\node[anchor=base,inner sep=0pt] at (1.600,-2.3) {$.$};
\node[anchor=base,inner sep=0pt] at (1.900,-2.3) {$0$};
\node[anchor=base,inner sep=0pt] at (2.700,-2.3) {$0$};
\node[anchor=base,inner sep=0pt] at (-.7,-1.1) {$+$};
\draw[line width=.5pt] (-1,-1.52)--(3.300,-1.52);
\node[anchor=base,inner sep=0pt,font=\fontsize{6.5}{7}\selectfont] at (1.900,.72) {$1$};
\node[anchor=base,inner sep=0pt,font=\fontsize{6.5}{7}\selectfont] at (0.800,.72) {$1$};
\node[anchor=base,inner sep=0pt,font=\fontsize{6.5}{7}\selectfont] at (0.000,.72) {$1$};
\end{tikzpicture}
\par\smallskip{\fontsize{10}{13}\selectfont 上方小字表示从后一位进来的数。\par}
\vspace{0pt}
\begin{enumerate}[leftmargin=1.7em,itemsep=5pt,topsep=0pt]
\item 百分位：\(\begin{aligned}[t]&\phantom{={}}8+2\\&=10\end{aligned}\)\quad 本位写0，向前一位进1。
\item 十分位：\(\begin{aligned}[t]&\phantom{={}}1+8+1\\&=9+1\\&=10\end{aligned}\)\quad 本位写0，向前一位进1。
\item 个位：\(\begin{aligned}[t]&\phantom{={}}7+7+1\\&=14+1\\&=15\end{aligned}\)\quad 本位写5，向前一位进1。
\end{enumerate}
最高位再写进上来的1。\par
\par\smallskip
因此，\(\begin{aligned}[t]&\phantom{={}}7.18+7.82\\&=15.00\end{aligned}\)。
\par\medskip

}

\QHeading{第3题　解决问题。}

\Problem{① 一本书的售价是18.95元，一支钢笔的售价是15.65元。李明带了40元，买一本书和一支钢笔，钱够吗？}
\par


\MethodTitle{方法一}
\Answer{{\Formula{\begin{aligned}[t]&\phantom{={}}40-(18.95+15.65)\\&=40-34.6\\&=5.4\text{（元）}\end{aligned}}}\par 答：钱够，还剩5.4元。}

\MethodTitle{方法二：分别买完后求剩余}
\Answer{先买书，再买钢笔，逐次求剩余的钱。}
\Answer{\Formula{\begin{aligned}[t]&\phantom{={}}40-18.95-15.65\\
&=21.05-15.65\\
&=5.4\text{（元）}\end{aligned}}}
\Answer{答：钱够，还剩5.4元。}
\Problem{② 一根跳绳长2.45米，另一根长1.8米，第一根比第二根长多少米？}

\MethodTitle{方法一}
\Answer{$2.45-1.8=0.65\text{（米）}$\par 答：第一根比第二根长0.65米。}

\MethodTitle{方法二：减2米，再加回多减的0.2米}
\Answer{1.8米比2米少0.2米。}
\Answer{\Formula{\begin{aligned}[t]&\phantom{={}}2.45-1.8\\
&=2.45-2+0.2\\
&=0.45+0.2\\
&=0.65\text{（米）}\end{aligned}}}
\Answer{答：第一根比第二根长0.65米。}

